\chapter{THIẾT KẾ, TRIỂN KHAI VÀ ĐÁNH GIÁ HỆ THỐNG}

Chương này trình bày cách UniTrack được tổ chức và triển khai, bao gồm kiến trúc tổng quan, route/API, mã nguồn, migration, kiểm thử, giao diện và mô hình triển khai tham khảo. Nội dung chương nhằm đối chiếu các yêu cầu đã phân tích với thiết kế và kết quả hiện thực hóa của hệ thống.

\section{Thiết kế kiến trúc tổng quan}

UniTrack sử dụng kiến trúc Web tách giao diện, API, cơ sở dữ liệu và file storage. Giao diện React chạy trong trình duyệt và gọi API bằng session cookie. Backend Go xử lý routing, strict JSON decoding, phân quyền, kiểm tra vòng đời dự án, transaction và kết nối storage. PostgreSQL lưu dữ liệu nghiệp vụ và các ràng buộc toàn vẹn. Minh chứng được lưu qua một lớp trừu tượng để môi trường phát triển cục bộ và Cloudflare R2 trong production dùng cùng một API.

\figurelead{fig:system-context.png}{đặt các lớp client, API, cơ sở dữ liệu, storage và các điểm kiểm soát request vào cùng một bức tranh kiến trúc để làm rõ ranh giới trách nhiệm}
\figSystemContext

Luồng yêu cầu điển hình gồm: trình duyệt gửi request kèm cookie; backend kiểm tra origin với yêu cầu ghi dữ liệu có session cookie hoặc route ghi session cookie; middleware xác thực session; handler kiểm tra route ID và nội dung request; helper phân quyền kiểm tra vai trò và quan hệ với dự án; các thao tác ghi khóa dòng dự án để kiểm tra lại quyền quản lý, vòng đời và đối tượng đích; ràng buộc cơ sở dữ liệu là lớp bảo vệ cuối cùng; response trả JSON hoặc stream file.

\figurelead{fig:request-boundary.png}{cụ thể hóa các lớp kiểm soát mà một request phải đi qua trước khi dữ liệu được đọc, ghi hoặc trả về cho người dùng}
\figRequestBoundary

\tableintro{tab:request-boundary-layers}{mô tả các lớp kiểm soát chính mà một request phải đi qua trước khi dữ liệu được ghi hoặc trả về}
\begin{table}[!htbp]
\caption{Các lớp kiểm soát ở request boundary}
\label{tab:request-boundary-layers}
\centering
\begin{tblr}{
  width=\textwidth,
  colspec={X[0.95,l]X[1.55,l]X[1.45,l]},
  rows={valign=t},
  row{1}={font=\bfseries,bg=black!7},
  hlines={0.35pt,black!24},
}
Lớp kiểm soát & Vai trò & Ví dụ áp dụng \\
Kiểm tra origin & Chặn yêu cầu ghi dữ liệu có session cookie hoặc route ghi session cookie khi origin không tin cậy & Đăng nhập, đăng xuất và các request ghi dữ liệu đã đăng nhập \\
Kiểm tra session & Đọc cookie, băm token, kiểm tra tài khoản active và session bị thu hồi & \route{GET /auth/me}, các route dự án cần đăng nhập \\
Strict JSON decoder & Giới hạn body, từ chối trường lạ và JSON nối tiếp & Tạo dự án, nhiệm vụ, duyệt bài, tài nguyên \\
Kiểm tra quan hệ & Xác định người dùng có quyền xem/quản lý dự án hay không & Chi tiết dự án, chi tiết nhiệm vụ, tải minh chứng \\
Khóa transaction & Đọc lại quyền quản lý, trạng thái dự án và đối tượng đích trong transaction trước khi ghi & Tạo nhiệm vụ, nộp bài, duyệt bài, tải minh chứng \\
Ràng buộc dữ liệu & Bảo vệ luật nghiệp vụ khi có request đồng thời hoặc ghi trực tiếp vào cơ sở dữ liệu & Một bài chờ duyệt, đích minh chứng cùng dự án, khóa minh chứng sau duyệt \\
\end{tblr}
\end{table}

\section{Thiết kế route và API}

\subsection{Route giao diện}

Giao diện sử dụng React Router. Các route chính bám sát sản phẩm: đăng nhập, dashboard, Workspace, chi tiết thư mục, chi tiết dự án, chi tiết nhiệm vụ và quản trị tài khoản. Các đường dẫn tương thích \route{/projects} và \route{/classes} được redirect về namespace \route{/workspace} để giữ hướng điều hướng project-first. Router tự guard đăng nhập, quản trị và chi tiết thư mục; quyền xem dự án/nhiệm vụ vẫn do API quyết định cuối cùng và giao diện hiển thị trạng thái bị từ chối khi API trả lỗi.

\tableintro{tab:frontend-routes}{tổng hợp các route giao diện đang triển khai và điều kiện truy cập tương ứng}
\begin{longtblr}[
  caption={Bảng route giao diện đang triển khai},
  label={tab:frontend-routes},
]{
  width=\textwidth,
  colspec={X[1.25,l]X[0.9,l]X[1.45,l]},
  rowhead=1,
  rows={valign=t},
  row{1}={font=\bfseries,bg=black!7},
  hlines={0.35pt,black!24},
}
Route giao diện & Điều kiện truy cập & Mục đích \\
\route{/login} & Công khai & Đăng nhập bằng email/mật khẩu \\
\route{/dashboard} & Đã đăng nhập & Hàng đợi công việc và tóm tắt theo vai trò \\
\route{/workspace} & Đã đăng nhập & Trang Workspace với thư mục và danh sách dự án \\
\route{/admin/users} & Quản trị viên & Quản lý tài khoản \\
\stackcell{\route{/workspace/classes/}\\\route{:classId}} & Giảng viên/quản trị viên & Chi tiết thư mục và thêm dự án có sẵn \\
\stackcell{\route{/workspace/projects/}\\\route{:projectId}} & Người có quyền xem dự án & Hồ sơ dự án, kế hoạch, nhóm và tài nguyên \\
\stackcell{\route{/workspace/projects/}\\\route{:projectId/tasks/}\\\route{:taskId}} & Người có quyền xem dự án & Chi tiết nhiệm vụ, nộp bài và duyệt bài \\
\end{longtblr}

\subsection{Route API backend}

Backend đăng ký các route nghiệp vụ dưới tiền tố \route{/api/v1}; phần lớn route được bảo vệ bởi \code{requireAuth}. Thay vì liệt kê toàn bộ endpoint, Bảng~\ref{tab:api-routes} gom API theo mục đích nghiệp vụ để người đọc thấy được dữ liệu nào được đọc/ghi và lớp kiểm soát nào đi kèm.

\tableintro{tab:api-routes}{tóm tắt các nhóm API backend theo mục đích nghiệp vụ và điều kiện kiểm soát chính}
\begin{longtblr}[
  caption={Các nhóm API backend chính của UniTrack},
  label={tab:api-routes},
]{
  width=\textwidth,
  colspec={X[0.9,l]X[1.45,l]X[1.45,l]},
  rowhead=1,
  rows={valign=t},
  row{1}={font=\bfseries,bg=black!7},
  hlines={0.35pt,black!24},
}
Nhóm API & Dữ liệu/thao tác chính & Điều kiện kiểm soát nổi bật \\
Auth & Đăng nhập, đăng xuất và đọc người dùng hiện tại & Kiểm tra origin, giới hạn đăng nhập, tài khoản active và session bị thu hồi \\
Quản trị tài khoản & Tạo tài khoản, tìm kiếm, đổi vai trò/trạng thái và đặt mật khẩu & Chỉ quản trị viên; không làm mất quản trị viên active cuối cùng; thu hồi session khi cần \\
Dashboard & Hàng đợi công việc theo vai trò & Chỉ trả dữ liệu trong phạm vi vai trò và quan hệ dự án của người dùng \\
Workspace và dự án & Danh sách thư mục, dự án, chi tiết dự án và gán dự án vào thư mục & Chủ thư mục phải khớp giảng viên phụ trách; dự án cũ phải được kiểm tra lại trước khi gán \\
Nhóm dự án & Danh sách thành viên, thêm/xóa sinh viên và đổi leader/member & Chỉ sinh viên active; trạng thái dự án phải cho phép thay đổi nhóm \\
Kế hoạch và nhiệm vụ & Mốc kế hoạch, thứ tự mốc, nhiệm vụ và người được giao & Dự án active khi tạo mới; mốc và sinh viên phải thuộc cùng dự án \\
Bài nộp và duyệt bài & Sinh viên nộp tiến độ; giảng viên/quản trị viên ghi quyết định duyệt & Một bài chờ duyệt cho mỗi nhiệm vụ; người duyệt phải còn quyền quản lý dự án \\
Tài nguyên và minh chứng & Đường dẫn tài nguyên, file minh chứng và tải xuống file riêng tư & Đích tài nguyên/minh chứng phải cùng dự án; dữ liệu của bài đã duyệt chuyển sang chỉ đọc \\
\end{longtblr}

\begin{artifactnote}{Request boundary}
Các JSON handler dùng strict decoder: giới hạn body 1 MB, từ chối trường lạ và từ chối JSON nối tiếp. Route ID của nhiệm vụ, mốc kế hoạch, bài nộp, thành viên, tài nguyên và file được kiểm tra UUID trước khi truy vấn để tránh lộ lỗi ép kiểu từ cơ sở dữ liệu ra client. Các lỗi validation mong đợi dùng status error có kiểu để không biến lỗi hạ tầng thành lỗi \code{400}.
\end{artifactnote}

Các danh sách có thể dài như dự án, tài khoản, thư mục, thành viên, nhiệm vụ và tài nguyên dùng phân trang để giao diện không phải tải quá nhiều dữ liệu một lần. Một số route lồng vẫn trả dạng mảng khi không truyền tham số phân trang để không làm đứt các màn hình đã có; khi có \code{page} hoặc \code{limit}, response trả thêm \code{items}, \code{page}, \code{limit} và \code{total}.

\section{Thiết kế gói mã nguồn}

Hai sơ đồ dưới đây tách riêng phụ thuộc gói của giao diện và backend để người đọc nhìn rõ từng phần triển khai trước khi chuyển sang thiết kế dữ liệu. Cách chia gói theo trách nhiệm cũng phù hợp với nguyên tắc module hóa, trong đó mỗi module nên che giấu chi tiết thay đổi của mình và chỉ công khai phần giao tiếp cần thiết \cite{parnas1972criteria}.

\subsection{Sơ đồ gói giao diện}

\figurelead{fig:frontend-package-map.png}{thể hiện các nhóm mã giao diện và quan hệ phụ thuộc giữa app shell, component dùng chung, feature page và lớp gọi API}
\figFrontendPackageMap

\subsection{Sơ đồ gói backend}

\figurelead{fig:backend-package-map.png}{thể hiện cách backend gom route, middleware, handler nghiệp vụ, cấu hình, storage và truy cập cơ sở dữ liệu}
\figBackendPackageMap

\subsection{Thiết kế thành phần chủ đạo}

Do backend Go không tổ chức theo lớp hướng đối tượng truyền thống, phần thiết kế chi tiết được trình bày theo các struct, component React và module nghiệp vụ chủ đạo. Các thành phần trong Bảng~\ref{tab:key-components} đại diện cho các điểm điều phối quan trọng nhất: khởi tạo route/API, xác thực, ghi dữ liệu theo vòng đời dự án, quản lý file minh chứng và hiển thị hồ sơ dự án trên giao diện.

\tableintro{tab:key-components}{mô tả các thành phần chủ đạo, trách nhiệm và dữ liệu/phương thức chính của từng thành phần}
\begingroup\small
\begin{longtblr}[
  caption={Thiết kế thành phần chủ đạo của UniTrack},
  label={tab:key-components},
]{
  width=\textwidth,
  colspec={X[0.85,l]X[1.35,l]X[1.45,l]X[1.35,l]},
  colsep=3pt,
  rowhead=1,
  rows={valign=t},
  row{1}={font=\bfseries,bg=black!7},
  hlines={0.35pt,black!24},
}
Thành phần & Trách nhiệm thiết kế & Cơ chế chính & Khu vực triển khai \\
Server & Gom cấu hình, pool cơ sở dữ liệu, logger, router, file store và rate limit store để đăng ký route/middleware thống nhất & Router, database pool, file store, login limiter & Backend API core \\
Nhóm auth và account & Điều phối đăng nhập, current user, logout, thao tác tài khoản và thu hồi session khi tài khoản thay đổi & Request DTO, session token hash, khóa dòng user, session revocation & Auth/account backend \\
Nhóm handler project/task & Bảo vệ luồng dự án, thành viên, mốc kế hoạch, nhiệm vụ, bài nộp và duyệt bài bằng kiểm tra quyền/vòng đời trong transaction & Helper phân quyền, \smallcode{SELECT ... FOR UPDATE}, trạng thái dự án/task/submission & Project/task backend \\
File store và triển khai & Trừu tượng hóa local filesystem và R2 để file minh chứng có cùng luồng upload, download và delete ở mọi môi trường & \smallcode{Put}, \smallcode{Open}, \smallcode{Delete}, metadata file, object key & Storage backend \\
Trang chi tiết dự án & Gom dữ liệu dự án, nhóm, mốc, nhiệm vụ và tài nguyên thành hồ sơ làm việc chính của giảng viên/sinh viên & Query keys, action guards, Work Plan, team popover, resource dialog & Project UI \\
Trang chi tiết nhiệm vụ & Hiển thị ngữ cảnh nhiệm vụ, lịch sử bài nộp, minh chứng, form nộp bài và vùng duyệt bài theo vai trò & Submission form, evidence panel, review form, stale refresh handlers & Assignment UI \\
\end{longtblr}
\endgroup

Thiết kế trên giữ ranh giới rõ giữa giao diện và backend. Giao diện chịu trách nhiệm trình bày trạng thái và hướng dẫn thao tác; backend quyết định quyền, vòng đời và tính hợp lệ dữ liệu; PostgreSQL là lớp cuối cùng bảo vệ invariant khi có request đồng thời hoặc ghi trực tiếp.

\section{Thiết kế dữ liệu}

Thiết kế dữ liệu đặt \code{projects} làm trục chính. Các bảng nhiệm vụ, bài nộp, duyệt bài, tài nguyên và file đều có cách liên kết để xác định dự án chứa chúng. Những quan hệ này phục vụ cả kiểm tra quyền lẫn ràng buộc toàn vẹn dữ liệu.

\subsection{Sơ đồ thực thể liên kết}

Sơ đồ thực thể liên kết tổng quan được đặt trên các trang riêng theo chiều dọc để giữ cùng hướng đọc với toàn bộ báo cáo. Hình~\ref{fig:data-model-core.png} dùng dạng thẻ bảng tương tự sơ đồ ERD tham khảo: header xanh đậm, thân bảng xám nhạt, field căn trái và kiểu dữ liệu căn phải. Sơ đồ liệt kê đầy đủ các field hiện có sau khi chạy toàn bộ migration \code{Up} và được tách thành ba phần liên tục để các bảng đủ lớn khi in trên khổ giấy dọc. Phần từ điển dữ liệu ngay sau sơ đồ bổ sung vai trò, ràng buộc và ghi chú nghiệp vụ cho từng bảng.

\clearpage
\begin{figure}[p]
  \centering
  \makebox[\textwidth][c]{\includegraphics[width=1.18\textwidth,height=0.84\textheight,keepaspectratio]{data-model-core.png}}
  \caption{ERD tổng quan schema dữ liệu và quan hệ chính của UniTrack}
  \label{fig:data-model-core.png}
  \caption*{\small Phần 1 trình bày nhóm tài khoản, phiên, Workspace, dự án và kế hoạch. Mỗi thẻ liệt kê field ở cột trái và kiểu dữ liệu ở cột phải; đường liền biểu diễn khóa ngoại chính, đường chấm biểu diễn quan hệ schema-only hoặc bổ trợ.}
\end{figure}
\clearpage
\begin{figure}[p]
  \ContinuedFloat
  \centering
  \makebox[\textwidth][c]{\includegraphics[width=1.18\textwidth,height=0.84\textheight,keepaspectratio]{data-model-core-support.png}}
  \caption[]{ERD tổng quan schema dữ liệu và quan hệ chính của UniTrack (tiếp theo)}
  \caption*{\small Phần 2 trình bày assignment, phân công, bài nộp và review. Các bảng ngữ cảnh như \code{users}, \code{projects} và \code{tasks} được lặp lại theo cùng định dạng để giữ quan hệ đọc được trên trang dọc.}
\end{figure}
\clearpage
\begin{figure}[p]
  \ContinuedFloat
  \centering
  \makebox[\textwidth][c]{\includegraphics[width=1.18\textwidth,height=0.84\textheight,keepaspectratio]{data-model-core-support-ops.png}}
  \caption[]{ERD tổng quan schema dữ liệu và quan hệ chính của UniTrack (tiếp theo)}
  \caption*{\small Phần 3 trình bày tài nguyên, minh chứng, cleanup và audit. Các bảng ngữ cảnh như \code{users}, \code{projects} và \code{progress\_updates} được lặp lại theo cùng định dạng để giữ quan hệ đọc được trên trang dọc.}
\end{figure}
\clearpage

\subsection{Chi tiết bảng dữ liệu}

Phần này trình bày từ điển dữ liệu ở trạng thái sau khi chạy toàn bộ migration \code{Up}. Các bảng đã bị loại khỏi schema hiện tại như \code{invitations}, \code{meeting\_notes}, \code{meeting\_note\_attendees}, \code{feedback} và \code{feedback\_replies} không được đưa vào từ điển. Hai bảng \code{courses} và \code{assessments} vẫn còn trong schema vì lịch sử migration; phần dưới đánh dấu rõ là schema-only để không nhầm với luồng chức năng hiện tại.

\begingroup\scriptsize
\tableintro{tab:data-dictionary-identity}{trình bày trường dữ liệu và ràng buộc chính của nhóm tài khoản, phiên và Workspace}
\begin{longtblr}[
  caption={Từ điển dữ liệu nhóm tài khoản, phiên và Workspace},
  label={tab:data-dictionary-identity},
]{
  width=\textwidth,
  colspec={X[0.95,l]X[1.55,l]X[1.70,l]},
  rowhead=1,
  rows={valign=t},
  row{1}={font=\bfseries,bg=black!7},
  hlines={0.35pt,black!24},
}
Bảng & Trường dữ liệu & Ràng buộc và vai trò \\
\code{users} & \code{id}; \code{full\_name}, \code{email}, \code{password\_hash}, \code{role}, \code{status}, \code{avatar\_url}; \code{created\_at}, \code{updated\_at}. & Khóa chính người dùng; email unique theo chữ thường; \code{role} và \code{status} là enum kiểm soát quyền đăng nhập và phân quyền. \\
\code{sessions} & \code{id}, \code{user\_id}; \code{token\_hash}, \code{expires\_at}, \code{revoked\_at}, \code{last\_seen\_at}, \code{user\_agent}, \code{ip\_address}, \code{created\_at}. & FK tới \code{users}; token chỉ lưu dạng hash; phiên active là phiên chưa thu hồi và chưa hết hạn. \\
\code{courses} & \code{id}; \code{code}, \code{title}, \code{description}, \code{status}; \code{created\_by}, \code{created\_at}, \code{updated\_at}. & Bảng schema-only lịch sử; \code{code} unique theo chữ thường; không còn route/giao diện course độc lập. \\
\code{course\_sections} & \code{id}; \code{title}, \code{description}, \code{color}, \code{status}; \code{owner\_teacher\_id}, \code{created\_by}, \code{created\_at}, \code{updated\_at}. & Thư mục Workspace; chủ thư mục phải là giảng viên; trigger bảo đảm thư mục chỉ chứa dự án do đúng giảng viên phụ trách. \\
\stackcell{\code{course\_section}\\\code{\_projects}} & \code{id}; \code{course\_section\_id}, \code{project\_id}, \code{added\_by}, \code{added\_at}. & Bảng nối thư mục--dự án; unique \code{project\_id} để một dự án chỉ nằm trong tối đa một thư mục. \\
\end{longtblr}

\tableintro{tab:data-dictionary-project}{trình bày trường dữ liệu và ràng buộc chính của nhóm dự án, nhóm thành viên, mốc kế hoạch và schema-only assessment}
\begin{longtblr}[
  caption={Từ điển dữ liệu nhóm dự án và kế hoạch},
  label={tab:data-dictionary-project},
]{
  width=\textwidth,
  colspec={X[0.95,l]X[1.55,l]X[1.70,l]},
  rowhead=1,
  rows={valign=t},
  row{1}={font=\bfseries,bg=black!7},
  hlines={0.35pt,black!24},
}
Bảng & Trường dữ liệu & Ràng buộc và vai trò \\
\code{projects} & \code{id}; \code{name}, \code{description}, \code{topic}; \code{supervisor\_id}, \code{created\_by}; \code{start\_date}, \code{end\_date}; \code{status}, \code{official\_progress\_state}, \code{progress\_summary}; \code{created\_at}, \code{updated\_at}. & Trục dữ liệu chính; check ngày kết thúc không trước ngày bắt đầu; \code{status} chi phối thao tác ghi theo vòng đời. \\
\code{project\_members} & \code{id}; \code{project\_id}, \code{student\_id}; \code{member\_role}; \code{joined\_at}. & Unique cặp dự án--sinh viên; partial unique bảo đảm mỗi dự án có tối đa một leader. \\
\code{project\_milestones} & \code{id}; \code{project\_id}, \code{created\_by}; \code{title}, \code{description}; \code{target\_date}, \code{sort\_order}; \code{created\_at}, \code{updated\_at}. & Mốc Work Plan; unique cặp \code{project\_id}, \code{id} để assignment và support kiểm tra cùng dự án. \\
\code{assessments} & \code{id}; \code{project\_id}, \code{assessed\_by}; \code{target\_type}, \code{target\_id}; \code{score}, \code{max\_score}, \code{note}; \code{created\_at}, \code{updated\_at}. & Bảng schema-only của luồng điểm số đã loại khỏi UI/API; vẫn có FK dự án/người chấm và check điểm không vượt thang điểm. \\
\end{longtblr}

\tableintro{tab:data-dictionary-assignment}{trình bày trường dữ liệu và ràng buộc chính của nhóm assignment, bài nộp và duyệt bài}
\begin{longtblr}[
  caption={Từ điển dữ liệu nhóm assignment và bài nộp},
  label={tab:data-dictionary-assignment},
]{
  width=\textwidth,
  colspec={X[0.95,l]X[1.55,l]X[1.70,l]},
  rowhead=1,
  rows={valign=t},
  row{1}={font=\bfseries,bg=black!7},
  hlines={0.35pt,black!24},
}
Bảng & Trường dữ liệu & Ràng buộc và vai trò \\
\code{tasks} & \code{id}; \code{project\_id}, \code{parent\_task\_id}, \code{milestone\_id}; \code{title}, \code{description}; \code{status}, \code{priority}, \code{deadline}, \code{official\_progress\_state}; \code{created\_by}, \code{updated\_by}, \code{created\_at}, \code{updated\_at}. & Assignment chính thức phải thuộc một mốc cùng dự án; child-task còn ở schema nhưng không còn workflow chủ động. \\
\code{task\_assignees} & \code{id}; \code{task\_id}, \code{project\_id}, \code{student\_id}; \code{assigned\_at}. & Unique cặp task--sinh viên; FK composite bảo đảm người được giao là thành viên cùng dự án. \\
\code{progress\_updates} & \code{id}; \code{project\_id}, \code{task\_id}, \code{submitted\_by}; \code{title}, \code{description}, \code{blockers}; \code{review\_status}; \code{created\_at}, \code{updated\_at}. & Một task chỉ có một bài \code{pending\_review}; trigger bảo đảm người nộp là sinh viên active được giao assignment chính thức. \\
\code{progress\_reviews} & \code{id}; \code{progress\_update\_id}, \code{reviewed\_by}; \code{review\_status}, \code{review\_comment}, \code{official\_progress\_state}; \code{reviewed\_at}. & Unique một quyết định cuối cùng cho mỗi bài nộp; backend kiểm tra người duyệt còn quyền quản lý dự án. \\
\end{longtblr}

\tableintro{tab:data-dictionary-support}{trình bày trường dữ liệu và ràng buộc chính của nhóm tài nguyên, minh chứng, cleanup và audit}
\begin{longtblr}[
  caption={Từ điển dữ liệu nhóm tài nguyên, minh chứng và audit},
  label={tab:data-dictionary-support},
]{
  width=\textwidth,
  colspec={X[0.95,l]X[1.55,l]X[1.70,l]},
  rowhead=1,
  rows={valign=t},
  row{1}={font=\bfseries,bg=black!7},
  hlines={0.35pt,black!24},
}
Bảng & Trường dữ liệu & Ràng buộc và vai trò \\
\code{resource\_links} & \code{id}; \code{project\_id}, \code{added\_by}; \code{related\_entity\_type}, \code{related\_entity\_id}; \code{title}, \code{url}, \code{type}, \code{description}; \code{created\_at}, \code{updated\_at}. & Liên kết đa hình tới project/milestone/task/submission; trigger kiểm tra đích cùng dự án; unique URL theo đích; khóa sau duyệt. \\
\code{uploaded\_files} & \code{id}; \code{project\_id}, \code{uploaded\_by}; \code{related\_entity\_type}, \code{related\_entity\_id}; \code{original\_file\_name}, \code{stored\_file\_name}, \code{storage\_path}, \code{mime\_type}, \code{file\_size\_bytes}, \code{created\_at}. & Metadata file riêng tư; trigger kiểm tra đích cùng dự án và khóa file gắn với bài đã duyệt. \\
\stackcell{\code{uploaded\_file}\\\code{\_object\_cleanup}\\\code{\_jobs}} & \code{id}; \code{storage\_path}; \code{reason}, \code{attempts}, \code{last\_error}; \code{completed\_at}, \code{created\_at}, \code{updated\_at}. & Hàng đợi dọn object khi upload/delete lệch giữa DB và storage; index pending theo job chưa có \code{completed\_at}. \\
\code{activity\_logs} & \code{id}; \code{actor\_id}, \code{project\_id}; \code{action}, \code{entity\_type}, \code{entity\_id}; \code{metadata}; \code{created\_at}. & Audit thao tác tài khoản và nhóm dự án; \code{project\_id} có thể rỗng với thao tác không thuộc dự án. \\
\end{longtblr}
\endgroup

\subsection{Migration và ràng buộc toàn vẹn nổi bật}

\tableintro{tab:key-migrations}{tóm tắt các nhóm migration quan trọng và mục đích nghiệp vụ mà từng nhóm bảo vệ}
\begin{longtblr}[
  caption={Các nhóm migration quan trọng và mục đích nghiệp vụ},
  label={tab:key-migrations},
]{
  width=\textwidth,
  colspec={X[0.95,l]X[1.25,l]X[1.55,l]},
  rowhead=1,
  rows={valign=t},
  row{1}={font=\bfseries,bg=black!7},
  hlines={0.35pt,black!24},
}
Nhóm migration & Khu vực & Vai trò trong thiết kế \\
Nền tảng dữ liệu & Tài khoản, session, dự án, nhiệm vụ và bài nộp & Tạo các bảng cốt lõi để hệ thống có đủ trục dữ liệu ban đầu. \\
Vòng đời và Workspace & Trạng thái dự án, thư mục Workspace và liên kết thư mục-dự án & Giữ hướng project-first, đồng thời bảo đảm dự án chỉ nằm trong thư mục phù hợp với giảng viên phụ trách. \\
Tài khoản và nhóm dự án & Email người dùng, session revocation, thành viên và leader & Bảo đảm email không trùng theo chữ thường, session bị thu hồi khi cần và mỗi dự án chỉ có một leader. \\
Kế hoạch và nhiệm vụ & Mốc kế hoạch, deadline theo ngày và nhiệm vụ bắt buộc thuộc mốc & Biến nhiệm vụ thành đơn vị công việc chính thức nằm trong một mốc kế hoạch rõ ràng. \\
Bài nộp và duyệt bài & Bài nộp cùng dự án, một bài chờ duyệt và người nộp hợp lệ & Chặn bài nộp lệch nhiệm vụ/dự án, chặn nhiều bài đang chờ cho cùng nhiệm vụ và bảo đảm người nộp được phân công. \\
Tài nguyên và minh chứng & Đích tài nguyên/file, khóa dữ liệu sau duyệt và file minh chứng riêng tư & Bảo đảm tài nguyên/minh chứng thuộc đúng dự án và không bị thay đổi sau khi đã là căn cứ duyệt bài. \\
Lưu trữ file & Hàng đợi dọn file lỗi & Giảm rủi ro lệch giữa metadata trong database và file thật khi upload/delete gặp lỗi giữa chừng. \\
\end{longtblr}

\figurelead{fig:project-lifecycle-full}{tóm tắt vòng đời dự án theo trạng thái, hướng chuyển trạng thái và mức độ ghi dữ liệu chính ở mỗi trạng thái}
\figProjectLifecycleFull{Vòng đời dự án và tác động tới thao tác ghi}

Hình~\ref{fig:project-lifecycle-full} đặt vòng đời dự án làm lớp kiểm soát bên ngoài vì mọi thao tác ghi thuộc Work Plan đều phải đi qua trạng thái dự án trước. Bên trong lớp này, nhiệm vụ và mốc kế hoạch vẫn có trạng thái riêng để biểu diễn tiến độ và kết quả duyệt. Bảng~\ref{tab:lifecycle-state-hierarchy} tóm tắt cách ba tầng trạng thái phối hợp với nhau.

\tableintro{tab:lifecycle-state-hierarchy}{phân biệt vai trò của vòng đời dự án, trạng thái nhiệm vụ và trạng thái mốc kế hoạch trong Work Plan}
\begin{table}[!htbp]
\caption{Phân tầng trạng thái trong Work Plan}
\label{tab:lifecycle-state-hierarchy}
\centering
\small
\begin{tblr}{
  width=\textwidth,
  colspec={X[0.72,l]X[1.40,l]X[1.68,l]},
  rows={valign=t},
  row{1}={font=\bfseries,bg=black!7},
  hlines={0.35pt,black!24},
}
Tầng trạng thái & Trạng thái chính & Vai trò trong kiểm soát nghiệp vụ \\
Dự án & \code{active}, \code{on\_hold}, \code{completed}, \code{archived} & Lớp bao ngoài quyết định dự án còn cho phép tạo/sửa nhóm, mốc, nhiệm vụ, tài nguyên và bài nộp hay chỉ còn đọc/xử lý phần tồn đọng. \\
Nhiệm vụ & \code{todo}, \code{in\_progress}, \code{submitted}, \code{needs\_changes}, \code{done}; trạng thái chính thức gồm \code{no\_progress}, \code{in\_progress}, \code{needs\_changes}, \code{completed} & Đơn vị công việc chính thức của sinh viên; bài nộp và quyết định duyệt cập nhật tiến độ, chặn trạng thái mâu thuẫn và đóng nhiệm vụ khi đã hoàn thành. \\
Mốc kế hoạch & \code{empty}, \code{planned}, \code{in\_progress}, \code{needs\_changes}, \code{completed} & Trạng thái dẫn xuất từ các nhiệm vụ con, giúp người quản lý nhìn nhanh sức khỏe của một giai đoạn mà không tạo thêm luồng duyệt riêng cho mốc. \\
\end{tblr}
\end{table}
\FloatBarrier

\section{Thiết kế luồng nghiệp vụ chi tiết}

\subsection*{Trình tự nộp bài, minh chứng và duyệt bài}

\figurelead{fig:evidence-sequence.png}{đặt các bước nộp bài, tải minh chứng và duyệt bài vào cùng một trình tự để làm rõ trách nhiệm của giao diện, backend, storage và cơ sở dữ liệu}
\figEvidenceSequence

\tableintro{tab:sequence-coverage}{mô tả phạm vi biểu đồ trình tự và các luồng tương đương được đặc tả bằng tương tác giữa giao diện, backend và cơ sở dữ liệu}
\begin{longtblr}[
  caption={Phạm vi thiết kế trình tự cho các use case chính},
  label={tab:sequence-coverage},
]{
  width=\textwidth,
  colspec={X[0.9,l]X[1.45,l]X[1.55,l]},
  rowhead=1,
  rows={valign=t},
  row{1}={font=\bfseries,bg=black!7},
  hlines={0.35pt,black!24},
}
Use case & Đối tượng tham gia & Trình tự xử lý chính \\
Đăng nhập & Browser, auth API, bảng \code{users}, bảng \code{sessions}, session cookie & Browser gửi email/mật khẩu; API chuẩn hóa email, áp dụng rate limit, khóa dòng user, kiểm tra mật khẩu, tạo token, lưu token hash và đặt cookie HttpOnly. \\
Tạo nhiệm vụ & Project detail page, task API, bảng \code{projects}, \code{project\_milestones}, \code{task\_assignees} & Giao diện gửi form; API khóa dự án, kiểm tra lại quyền quản lý/vòng đời/mốc kế hoạch/sinh viên được giao, ghi task và assignee trong transaction, sau đó refresh project/task/dashboard. \\
Nộp bài, tải minh chứng và duyệt bài & Task detail page, progress API, file API, storage, bảng \code{progress\_updates}, \code{uploaded\_files}, hàng đợi cleanup object và \code{progress\_reviews} & Luồng được minh họa ở Hình~\ref{fig:evidence-sequence.png}: sinh viên nộp bài, upload minh chứng, giảng viên duyệt, sau đó support record chuyển sang chỉ đọc. \\
\end{longtblr}

\tableintro{tab:evidence-sequence-responsibilities}{phân rã trách nhiệm của giao diện, backend và cơ sở dữ liệu trong từng giai đoạn của luồng nộp bài, minh chứng và duyệt bài}
\begin{table}[H]
\caption{Trách nhiệm từng lớp trong luồng nộp bài, minh chứng và duyệt bài}
\label{tab:evidence-sequence-responsibilities}
\centering
\begin{tblr}{
  width=\textwidth,
  colspec={X[0.8,l]X[1.55,l]X[1.65,l]},
  rows={valign=t},
  row{1}={font=\bfseries,bg=black!7},
  hlines={0.35pt,black!24},
}
Giai đoạn & Trách nhiệm giao diện & Trách nhiệm backend/cơ sở dữ liệu \\
Mở nhiệm vụ & Tải chi tiết nhiệm vụ, dự án, thành viên, mốc kế hoạch, tài nguyên và file; chặn duyệt bài nếu dữ liệu hỗ trợ chưa sẵn sàng & Trả dữ liệu theo quyền xem dự án và không lộ dự án ngoài phạm vi \\
Nộp bài & Kiểm tra form, gửi tiêu đề/mô tả/vướng mắc và refresh dữ liệu liên quan sau khi nộp & Khóa dự án, kiểm tra lại thành viên/người được giao, chặn bài chờ duyệt trùng và ghi bài nộp \\
Tải minh chứng & Kiểm tra dung lượng 10 MB, gửi file tới bài nộp đang chờ duyệt & Tiền kiểm thông tin file/đích/vòng đời trước khi ghi file, tạo job dọn lỗi, lưu file riêng tư, khóa dự án kiểm tra lại rồi ghi metadata; nếu metadata/commit lỗi thì job còn pending để dọn file \\
Duyệt bài & Hiển thị vùng duyệt bài, lựa chọn quyết định và yêu cầu hướng dẫn khi trả bài/từ chối & Khóa bài chờ duyệt, ghi quyết định cuối cùng, đồng bộ trạng thái nhiệm vụ và bảo đảm mỗi bài chỉ được duyệt một lần \\
Sau duyệt & Hiển thị lịch sử và minh chứng ở chế độ chỉ đọc & Trigger cơ sở dữ liệu chặn thêm/sửa/xóa tài nguyên hoặc minh chứng của bài đã duyệt \\
\end{tblr}
\end{table}

\section{Thiết kế giao diện}

Giao diện UniTrack dùng sắc xanh học thuật, thanh điều hướng đậm, các khối nội dung sáng như giấy, dấu trạng thái rõ và danh sách hành động trực tiếp. Phần ảnh minh họa không nhằm liệt kê mọi màn hình, mà chọn các điểm giúp người đọc hiểu luồng sản phẩm: người dùng đăng nhập, nhìn thấy việc cần xử lý, đi vào Workspace, mở hồ sơ dự án, thao tác với nhóm/kế hoạch/tài nguyên, rồi nộp hoặc duyệt bài. Bảng~\ref{tab:ui-screenshot-rationale} sau các hình tóm tắt mỗi ảnh đang chứng minh ý đồ thiết kế nào.

Ba ảnh đầu cho thấy cách người dùng bắt đầu công việc: đăng nhập, xem dashboard theo vai trò và đi vào Workspace để chọn dự án hoặc thư mục.

\screenshotfigure{login.png}{Giao diện đăng nhập UniTrack}

\screenshotfigure{dashboard-admin.png}{Dashboard quản trị với bài chờ duyệt, nhiệm vụ quá hạn và phần tóm tắt}

\screenshotfigure{workspace.png}{Workspace với thư mục dự án và danh sách dự án độc lập}

Nhóm ảnh tiếp theo tập trung vào hồ sơ dự án. Đây là nơi giảng viên đọc tình hình dự án, xem nhóm sinh viên, quản lý mốc kế hoạch và gắn tài nguyên đúng đối tượng nghiệp vụ.

\screenshotfigure{folder-detail.png}{Màn hình chi tiết thư mục và danh sách dự án trong thư mục}

\screenshotfigure{project-detail.png}{Chi tiết dự án với tiêu đề thao tác, nhóm và kế hoạch công việc}

\screenshotfigure{project-team-popover.png}{Popover nhóm dự án với giảng viên phụ trách, leader và danh sách sinh viên}

\screenshotfigure{project-manage-plan.png}{Menu thao tác mốc kế hoạch với sửa, sắp xếp và tài nguyên}

\screenshotfigure{project-resource-dialog.png}{Hộp thoại quản lý tài nguyên theo mốc kế hoạch trong dự án}

Ba ảnh cuối thể hiện hai đầu của vòng lặp nộp và duyệt bài: giảng viên có vùng quyết định kèm minh chứng, sinh viên có hộp thoại nộp bài, còn quản trị viên có màn hình tài khoản để kiểm soát người dùng tham gia hệ thống.

\screenshotfigure{assignment-review-desk.png}{Trang duyệt nhiệm vụ với minh chứng và vùng quyết định}

\screenshotfigure{assignment-resource-dialog.png}{Hộp thoại tài nguyên gắn với bài nộp đang chờ duyệt}

\screenshotfigure{assignment-submit-work.png}{Giao diện sinh viên mở hộp thoại nộp bài cho nhiệm vụ được giao}

\screenshotfigure{admin-accounts.png}{Giao diện quản trị tài khoản}

\FloatBarrier

\begingroup\small
\begin{table}[!htbp]
\caption{Đối chiếu ảnh giao diện với ý đồ thiết kế}
\label{tab:ui-screenshot-rationale}
\centering
\begin{tblr}{
  width=\textwidth,
  colspec={X[0.95,l]X[1.45,l]X[1.55,l]},
  rowsep=2pt,
  rows={valign=t},
  row{1}={font=\bfseries,bg=black!7},
  hlines={0.35pt,black!24},
}
Màn hình & Ý đồ thiết kế & Thành phần kiểm chứng được trong ảnh \\
Đăng nhập & Tập trung vào đăng nhập, không phân tán bằng dashboard hoặc nội dung quảng bá & Khung đăng nhập ở giữa, label rõ, lỗi thân thiện và nền xanh đậm \\
Dashboard & Đưa việc cần xử lý lên trước các số liệu tổng quát & Tóm tắt nhanh, bài chờ duyệt cũ nhất trước, danh sách quá hạn và dự án cần chú ý \\
Workspace & Tổ chức nhiều dự án bằng thư mục nhẹ và danh sách dự án & Kệ thư mục, tìm kiếm, trạng thái/cảnh báo và dự án độc lập \\
Chi tiết dự án & Biến dự án thành hồ sơ làm việc hàng ngày & Tiêu đề thao tác, dấu trạng thái/tiến độ, nhóm và kế hoạch theo mốc \\
Popover nhóm & Cho phép xem nhanh và quản lý thành viên mà không rời hồ sơ dự án & Giảng viên phụ trách, leader, sinh viên, thêm thành viên và thao tác vai trò \\
Thao tác mốc kế hoạch & Giữ kế hoạch dễ đọc nhưng vẫn cho phép sửa từng mốc khi cần & Nút Actions theo mốc, thêm nhiệm vụ, sắp xếp, chỉnh sửa mốc và mở tài nguyên \\
Hộp thoại tài nguyên & Gom link tài liệu theo đúng đích nghiệp vụ & Tài nguyên của mốc/nhiệm vụ/bài nộp, form thêm link và trạng thái trống \\
Trang duyệt nhiệm vụ & Giúp người duyệt quyết định nhanh nhưng vẫn thấy ngữ cảnh và minh chứng & Bài nộp, tài nguyên, file minh chứng và vùng quyết định \\
Nộp bài của sinh viên & Tập trung vào mô tả kết quả và vướng mắc trước khi gắn minh chứng & Hộp thoại nộp bài, tiêu đề, mô tả, blockers và hành động gửi \\
Quản trị tài khoản & Quản lý tài khoản không cần đăng ký công khai & Tìm kiếm/lọc, vai trò/trạng thái, đặt mật khẩu, chỉnh sửa và phân trang \\
\end{tblr}
\end{table}
\endgroup

\section{Công cụ phát triển và thống kê mã nguồn}

Phần này bổ sung các thông tin triển khai theo yêu cầu của báo cáo ứng dụng: công cụ, thư viện, phiên bản chính, địa chỉ tham khảo và quy mô mã nguồn. Các phiên bản trong Bảng~\ref{tab:tool-versions} được ghi nhận theo cấu hình dự án tại thời điểm viết báo cáo.

\tableintro{tab:tool-versions}{liệt kê công cụ và thư viện chính kèm phiên bản hoặc nguồn tham khảo chính thức}
\begin{longtblr}[
  caption={Danh sách công cụ, thư viện và phiên bản chính},
  label={tab:tool-versions},
]{
  width=\textwidth,
  colspec={X[0.8,l]X[0.9,l]X[0.95,l]X[1.65,l]},
  rowhead=1,
  rows={valign=t},
  row{1}={font=\bfseries,bg=black!7},
  hlines={0.35pt,black!24},
}
Mục đích & Công cụ/thư viện & Phiên bản & URL tham khảo \\
Ngôn ngữ backend & Go & 1.25.5 & \url{https://go.dev/} \\
Router backend & chi & v5.2.5 & \url{https://github.com/go-chi/chi} \\
Driver PostgreSQL & pgx & v5.9.2 & \url{https://github.com/jackc/pgx} \\
Storage R2/S3 & AWS SDK Go v2 & v1.42.0, S3 v1.104.0 & \url{https://aws.github.io/aws-sdk-go-v2/} \\
Giao diện & React, React DOM & 19.2.5 & \url{https://react.dev/} \\
Build frontend & Vite & 8.0.9 & \url{https://vite.dev/} \\
Ngôn ngữ frontend & TypeScript & 6.0.2 & \url{https://www.typescriptlang.org/} \\
State server & TanStack Query & 5.99.2 & \url{https://tanstack.com/query/latest} \\
E2E/browser test & Playwright & 1.61.0 & \url{https://playwright.dev/} \\
CSS utility & Tailwind CSS & 4.2.4 & \url{https://tailwindcss.com/} \\
\end{longtblr}

\tableintro{tab:source-statistics}{tóm tắt quy mô mã nguồn và các thành phần triển khai tại thời điểm viết báo cáo}
\begin{table}[H]
\caption{Thống kê mã nguồn và thành phần triển khai chính}
\label{tab:source-statistics}
\centering
\begin{tblr}{
  width=\textwidth,
  colspec={X[1.1,l]X[0.7,c]X[0.85,c]X[1.6,l]},
  rows={valign=t},
  row{1}={font=\bfseries,bg=black!7},
  hlines={0.35pt,black!24},
}
Khu vực & Số file & Số dòng & Ghi chú \\
Backend Go & 26 & 14.624 & Bao gồm API server, handler nghiệp vụ, config, storage và kiểm thử backend \\
Frontend TypeScript/TSX & 77 & 9.039 & Bao gồm app shell, component dùng chung, feature pages, API client và types \\
Tổng Go/TS/TSX chính & 103 & 23.663 & Chưa tính SQL migration, tài liệu phụ trợ và asset tĩnh \\
Database migration & 23 & -- & Schema nền tảng, ràng buộc toàn vẹn và hàng đợi dọn file lỗi \\
Kịch bản Playwright & 9 & -- & Browser flow cho đăng nhập, dashboard, phân quyền, accessibility, luồng assignment happy path và trạng thái giao diện đã lỗi thời \\
\end{tblr}
\end{table}

\section{Kết quả triển khai chức năng}

\tableintro{tab:implementation-results}{tổng hợp trạng thái triển khai các nhóm chức năng chính của UniTrack}
\begin{longtblr}[
  caption={Tổng hợp kết quả triển khai chức năng},
  label={tab:implementation-results},
]{
  width=\textwidth,
  colspec={X[0.85,l]X[1.95,l]X[1.0,l]},
  rowhead=1,
  rowsep=2pt,
  rows={valign=t},
  row{1}={font=\bfseries,bg=black!7},
  hlines={0.35pt,black!24},
}
Nhóm chức năng & Kết quả triển khai & Trạng thái \\
Xác thực/session & Đăng nhập/đăng xuất/current user bằng cookie, kiểm tra origin, giới hạn đăng nhập có xét trusted proxy khi cấu hình, parse cấu hình nghiêm ngặt và kiểm tra cookie/CORS lúc khởi động & \okbadge \\
Quản lý tài khoản & Tạo/tìm kiếm/cập nhật vai trò/trạng thái/mật khẩu, bảo vệ quản trị viên cuối cùng và ghi nhật ký thao tác & \okbadge \\
Dashboard & Hàng đợi theo vai trò, nhánh quản trị/giảng viên/sinh viên, bài chờ duyệt cũ nhất trước và nhiệm vụ sinh viên cần xử lý & \okbadge \\
Workspace/thư mục & Kệ thư mục màu, tìm kiếm, danh sách dự án độc lập, thêm dự án có sẵn an toàn với dữ liệu đã tải trước và khớp chủ sở hữu & \okbadge \\
Dự án/nhóm & Tạo/sửa/xem dự án, kiểm tra lại quyền quản lý và vòng đời trong transaction, tổng hợp tiến độ, kế hoạch công việc, nhóm và vai trò leader/member & \okbadge \\
\stackcell{Nhiệm vụ\\bài nộp\\duyệt bài} & Nhiệm vụ thuộc mốc kế hoạch, sinh viên được giao active, vùng nộp bài, vùng duyệt bài và một quyết định cuối cùng & \okbadge \\
Tài nguyên/minh chứng & Hộp thoại tài nguyên, tải lên/tải xuống/xóa minh chứng, tiền kiểm upload, hàng đợi dọn file lỗi khi khởi động, adapter R2 và khóa dữ liệu sau duyệt & \okbadge \\
Nhật ký hoạt động & Đã ghi thao tác tài khoản và thêm/đổi vai trò/xóa thành viên; chưa có giao diện tra cứu audit & \partialbadge \\
Kiểm thử giao diện & Playwright bao phủ các luồng rủi ro chính, nhưng chưa bao phủ toàn bộ form và biến thể tài nguyên/minh chứng & \partialbadge \\
\end{longtblr}

Mã nguồn được tổ chức theo lát tính năng thay vì tách thành các module sản phẩm rời rạc. Các phần backend chính nằm quanh xác thực/tài khoản, dashboard, dự án-nhiệm vụ và tài nguyên-minh chứng; các phần giao diện chính nằm quanh app shell, Workspace, hồ sơ dự án và trang nhiệm vụ. Cách chia này giúp người đọc đối chiếu kết quả triển khai với các nhóm nghiệp vụ đã nêu ở Chương 2 mà không cần nhớ toàn bộ cấu trúc thư mục mã nguồn.

\section{Kiểm thử và đánh giá chất lượng}

Kiểm thử tập trung vào rủi ro nghiệp vụ, không dừng ở luồng thuận. Các test quan trọng xác nhận route cần đăng nhập được bảo vệ, vòng đời và quyền quản lý dự án được kiểm tra lại khi ghi, cơ sở dữ liệu từ chối dữ liệu lệch dự án, tài khoản inactive không giữ session, minh chứng đã duyệt bị khóa chỉnh sửa và giao diện không giữ thao tác đã lỗi thời.

\tableintro{tab:test-techniques}{phân loại kỹ thuật kiểm thử được dùng và mục tiêu chất lượng của từng nhóm kiểm thử}
\begin{longtblr}[
  caption={Kỹ thuật kiểm thử và mục tiêu kiểm soát chất lượng},
  label={tab:test-techniques},
]{
  width=\textwidth,
  colspec={X[0.95,l]X[1.45,l]X[1.55,l]},
  rowhead=1,
  rows={valign=t},
  row{1}={font=\bfseries,bg=black!7},
  hlines={0.35pt,black!24},
}
Kỹ thuật & Cách áp dụng trong UniTrack & Rủi ro được kiểm soát \\
Unit/integration test backend & Gọi handler và database test để kiểm tra session, permission, vòng đời, bài nộp, review, storage và trigger & Lỗi nghiệp vụ, dữ liệu sai dự án, session cũ, request đồng thời \\
Kiểm thử biên và validation & Gửi JSON sai, route ID sai, trạng thái không hợp lệ, file quá lớn hoặc thao tác không đúng vòng đời & API trả lỗi sai, lộ lỗi database cast, ghi dữ liệu không hợp lệ \\
Kiểm thử đồng thời và regression & Mô phỏng đăng nhập đua với vô hiệu hóa, tạo nhiệm vụ đua với completed, duyệt bài cũ & Race condition ở quyền/vòng đời và account-control \\
Database integrity test & Ghi trực tiếp vào database để xác nhận FK, unique index và trigger vẫn chặn invariant bị phá & Lỗi bypass API hoặc migration làm lệch dữ liệu \\
Playwright E2E/accessibility test & Chạy trình duyệt qua login, dashboard, admin, access-control, assignment happy path, trạng thái giao diện lỗi thời, dialog và keyboard flow & Route guard sai, affordance lỗi thời, luồng chính bị đứt, focus/keyboard không đạt \\
\end{longtblr}

\tableintro{tab:representative-tests}{nêu các kịch bản kiểm thử tiêu biểu và ý nghĩa của từng kịch bản đối với chất lượng hệ thống}
\begin{longtblr}[
  caption={Các kịch bản kiểm thử tiêu biểu và ý nghĩa},
  label={tab:representative-tests},
]{
  width=\textwidth,
  colspec={X[1.05,l]X[1.35,l]X[1.35,l]},
  rowhead=1,
  rows={valign=t},
  row{1}={font=\bfseries,bg=black!7},
  hlines={0.35pt,black!24},
}
Kịch bản & Lớp kiểm thử & Ý nghĩa đối với UniTrack \\
Đăng nhập, đăng xuất, vô hiệu hóa tài khoản và đặt lại mật khẩu & Go test & Bảo đảm session cũ không tiếp tục hợp lệ khi tài khoản thay đổi. \\
Quản trị viên active cuối cùng & Go test & Tránh mất đường quản trị hệ thống do thao tác cập nhật sai hoặc đồng thời. \\
Tạo nhiệm vụ hoặc nộp bài khi dự án đổi trạng thái & Go test với database & Chứng minh server kiểm tra lại vòng đời dự án trước khi ghi dữ liệu. \\
Ghi trực tiếp dữ liệu lệch dự án & Database integrity test & Chứng minh khóa ngoại, unique index và trigger vẫn bảo vệ dữ liệu dù bỏ qua API. \\
Minh chứng sau duyệt & Go test và Playwright & Xác nhận file/tài nguyên đã dùng để duyệt bài chỉ còn đọc/tải xuống, không bị sửa hoặc xóa. \\
Assignment happy path & Playwright & Kiểm tra luồng người dùng đầy đủ: giảng viên tạo dự án, thêm sinh viên, giao nhiệm vụ; sinh viên nộp bài; giảng viên duyệt bài. \\
Access control và accessibility & Playwright & Kiểm tra route bị giới hạn đúng vai trò, người ngoài dự án bị chặn, hộp thoại/focus/keyboard vẫn dùng được. \\
\end{longtblr}

\tableintro{tab:verification-commands}{liệt kê các lệnh kiểm thử và build dùng để xác nhận chất lượng trước khi bàn giao}
\begin{table}[H]
\caption{Các lệnh kiểm thử và build chính}
\label{tab:verification-commands}
\centering
\begin{tblr}{
  width=\textwidth,
  colspec={X[1.15,l]X[2.2,l]},
  rows={valign=t},
  row{1}={font=\bfseries,bg=black!7},
  hlines={0.35pt,black!24},
}
Lệnh & Mục đích \\
\stackcell{\smallcode{TEST\_DATABASE\_URL=}\\\smallcode{'<postgres-url>'}\\\smallcode{make api-test}} & Chạy kiểm thử backend có database; lệnh bắt buộc cấu hình \code{TEST\_DATABASE\_URL} để test vòng đời không bị bỏ qua ngầm \\
\smallcode{make api-test-unit} & Chạy kiểm thử Go không bắt buộc database; các test DB tự bỏ qua khi thiếu \code{TEST\_DATABASE\_URL} \\
\code{make api-build} & Biên dịch API server \\
\code{make db-validate} & Kiểm tra migration bằng goose \\
\code{pnpm --filter @unitrack/web lint} & Kiểm tra lint/type của giao diện \\
\code{pnpm --filter @unitrack/web build} & Tạo bản build production cho giao diện \\
\code{pnpm --filter @unitrack/web test:e2e} & Chạy Playwright khi web/API khả dụng \\
\end{tblr}
\end{table}

\section{Triển khai tham khảo}

Kiến trúc triển khai đề xuất dùng Vercel cho giao diện, Render cho Go API, Neon cho PostgreSQL và Cloudflare R2 cho kho minh chứng. Hình~\ref{fig:deployment.png} cũng thể hiện các điểm vận hành cần cấu hình khi đưa hệ thống lên production: biến môi trường bắt buộc, health/readiness check, bước migration database và đường đi riêng tư của file minh chứng. Giai đoạn đầu chưa cần CDN hoặc load balancer riêng vì Vercel đã phục vụ giao diện tĩnh qua CDN và Render cung cấp HTTPS routing cho API.

\figurelead{fig:deployment.png}{mô tả phương án triển khai tham khảo và các thành phần hạ tầng cần cấu hình khi đưa hệ thống lên môi trường production}
\figDeployment

\tableintro{tab:deployment-config}{tóm tắt các thành phần triển khai đề xuất và cấu hình quan trọng cần thiết cho từng thành phần}
\begin{table}[H]
\caption{Cấu hình triển khai tham khảo}
\label{tab:deployment-config}
\centering
\begin{tblr}{
  width=\textwidth,
  colspec={X[0.85,l]X[0.9,l]X[2.15,l]},
  rows={valign=t},
  row{1}={font=\bfseries,bg=black!7},
  hlines={0.35pt,black!24},
}
Thành phần & Nền tảng & Cấu hình quan trọng \\
Giao diện & Vercel & \code{VITE\_API\_URL} trỏ tới Render API; SPA rewrite để refresh route không 404 \\
Backend & Render & \code{DATABASE\_URL}, \code{CORS\_ALLOWED\_ORIGINS} chính xác, \code{TRUSTED\_PROXY\_CIDRS} nếu tin header proxy, \code{SESSION\_SECURE=true}, \code{SESSION\_SAME\_SITE=none}; production cần bootstrap admin mạnh hoặc admin active sẵn \\
Cơ sở dữ liệu & Neon & PostgreSQL quản lý sẵn, chạy migration trước khi seed/sử dụng \\
Lưu trữ & Cloudflare R2 & \code{UPLOAD\_STORAGE\_BACKEND=r2}, bucket private, access key/endpoint cấu hình trong Render; download qua API và mỗi lần API khởi động xử lý tối đa 100 job dọn file còn pending \\
\end{tblr}
\end{table}

\section{Đánh giá hệ thống}

UniTrack đã đáp ứng quy trình chính của bài toán: tạo dự án, tổ chức Workspace, quản lý nhóm, tạo nhiệm vụ, sinh viên nộp bài, giảng viên duyệt bài và bảo toàn minh chứng. Điểm mạnh của hệ thống nằm ở các lớp bảo vệ cho quyền, vòng đời dự án và tính nhất quán dữ liệu ở backend/cơ sở dữ liệu, thay cho cách triển khai CRUD thuần túy. Điểm còn thiếu là quản lý toàn bộ dự án cho quản trị viên, giao diện audit log, hardening kho minh chứng trong production và độ phủ kiểm thử giao diện rộng hơn cho các form/tài nguyên/minh chứng.

\tableintro{tab:system-evaluation}{đánh giá tổng quan hệ thống theo nghiệp vụ, an toàn truy cập, toàn vẹn dữ liệu, bảo trì, kiểm thử và hạn chế}
\begin{table}[H]
\caption{Đánh giá tổng quan hệ thống}
\label{tab:system-evaluation}
\centering
\begin{tblr}{
  width=\textwidth,
  colspec={X[1.0,l]X[2.4,l]},
  rows={valign=t},
  row{1}={font=\bfseries,bg=black!7},
  hlines={0.35pt,black!24},
}
Tiêu chí & Đánh giá \\
Đáp ứng nghiệp vụ & Bao phủ luồng project-first từ Workspace, nhiệm vụ, bài nộp tới duyệt bài \\
An toàn truy cập & Có guard ở giao diện, kiểm tra quyền ở backend, kiểm tra tài khoản active và kiểm tra origin \\
Toàn vẹn dữ liệu & Có khóa transaction, khóa ngoại/trigger theo dự án và khóa tài nguyên/minh chứng sau duyệt \\
Khả năng bảo trì & Mã nguồn tổ chức theo lát tính năng; phần mô tả kỹ thuật bám sát triển khai \\
Kiểm thử & Backend bao phủ tốt vòng đời/toàn vẹn dữ liệu; Playwright tập trung vào rủi ro giao diện quan trọng \\
Hạn chế & Chưa có giao diện audit, quản lý dự án nâng cao cho quản trị viên, quota/MIME/malware scanning cho storage và kiểm thử giao diện rộng hơn \\
\end{tblr}
\end{table}

Phần thiết kế đã trình bày kiến trúc, route/API, gói mã nguồn, thành phần chủ đạo, dữ liệu, giao diện, công cụ triển khai, kết quả chức năng, kiểm thử và mô hình triển khai tham khảo. Các điểm có tính đóng góp cao hơn, đặc biệt là bảo mật session, khóa vòng đời, ràng buộc database và bảo toàn minh chứng, được phân tích sâu hơn ở Chương 5.
